% !TEX root = TI5.tex
\section{Preliminaries}

\subsection{Rings and modules}

We review some basic notions of rings and modules that are used in this article. 

Let $\rng$ be a commutative unital (=has a multiplicative identity) ring. Let $0$ and $1$ denote the additive and multiplicative identities, respectively. For a non-zero $r\in \rng$, $r$ is a unit if there exists $r'$ such that $rr'=1$. It is a zero-divisor if there exists $r'$ such that $rr'=0$. It is nilpotent if there exists $n\in \N$ such that $r^n=0$. A ring is a domain if it does not contain zero-divisors.


A (left-)module $M$ over $\rng$ is an abelian group with a multiplication from $\rng$, denoted as $\cdot:\rng\times M\to M$, satisfying certain natural conditions. A submodule of $M$ is a subgroup of $M$ closed under the multiplication of $\rng$. For $S\subseteq M$, the module generated by $S$ is the smallest submodule of $M$ containing $S$, denoted by $\langle S\rangle$. We say that $S$ is a generating set of $M$ if $\langle S\rangle=M$. Note that $\rng$ is naturally an $\rng$-module.

A free module is any module isomorphic to $\rng^n$ for some $n$, where $\rng^n$ consists of length-$n$ row vectors over $\rng$, and the multiplication by $\rng$ is the usual ``scalar-vector coordinate-wise'' multiplication $r \cdot (r_1, \dotsc, r_n) = (rr_1, \dotsc, rr_n)$. 

%is a module with a linearly independent basis. A typical example is $\rng^n$ which consists of length-$n$
Tensor products, symmetric powers, and exterior powers of modules over commutative rings are classical \cite{AM94}. Let $\rng$ be a commutative ring. Let $M$ and $N$ be $\rng$-modules. We use $M\otimes N$, $\S^k(M)$, and $\Lambda^k(M)$ to denote the tensor product of $M$ and $N$, the $k$th symmetric power of $M$, and the $k$th exterior power of $M$.

A subring $S$ of $\rng$ is a subset of $\rng$ closed under addition and multiplication. An ideal $I$ of $\rng$ is a subring that is also closed under multiplication from $\rng$, that is, it is a submodule of $\rng$ as an $\rng$-module. Given an ideal $I$, the quotient ring $\rng/I$ is naturally defined. We say that $\rng$ is local, if it has a unique maximal ideal. An ideal $I$ of $\rng$ is principal if it can be generated by a single element. We say that $\rng$ is a principal ideal ring (PIR) if every ideal is principal. 

The annihilator of $r \in \rng$ is $\Ann_\rng(r) = \{s \in \rng : sr = 0\}$. For a subset $T \subseteq \rng$, its annihilator is $\Ann_\rng(T) = \bigcap_{r \in T} \Ann_\rng(r) = \{s \in \rng : st = 0 \forall t \in T\}$. For $r \in \rng, T \subseteq \rng$, $\Ann_\rng(r)$ and $\Ann_\rng(T)$ are both ideals. Sometimes we may write just $\Ann(r), \Ann(T)$, if $\rng$ is clear from context.

\subsection{Matrices and tensors over commutative rings}

In this article our base ring is commutative with a multiplicative identity. 

Let $\rng$ be such a ring. The set of $m\times n$ matrices over $\rng$ is denoted by $\M(m\times n, \rng)$. When $n=m$, we may write $\M(n, \rng)$ for $\M(m\times n, \rng)$. Let $\idmat_n$ be the $n\times n$ identity matrix in $\M(n, \rng)$. 

We say that $A\in\M(n, \rng)$ is \emph{invertible}, if there exists $A'\in \M(n, \rng)$, such that $AA'=A'A=\idmat_n$. Let $\GL(n, \rng)$ be the group of $n\times n$ invertible matrices over $\rng$. 

For $A\in \M(n, \rng)$, the determinant of $A$ is denoted by $\det(A)$. It is a basic fact that $A\in \GL(n, \rng)$ if and only if $\det(A)$ is a unit in $\rng$. 

The set of length-$n$ row vectors over $\rng$ is denoted by $\rng^n$. The set of length-$n$ column vectors over $\rng$ is denoted by $^n\rng$.

For $A\in \M(m\times n, \rng)$, the row module is the submodule of $\rng^n$ consisting of $\rng$ combinations of rows of $A$. The column module is defined similarly. The left kernel of $A$ is the submodule of $\rng^m$ consisting of $v\in\rng^m$ such that $vA=0$. The right kernel of $A$ is the submodule of $^n\rng$ consisting of $w\in{}^n\rng$ such that $Aw=0$.

\paragraph{Block matrices over commutative rings.} 
For two matrices $A \in \M(m_1 \times n_1, \rng), B \in \M(m_2 \times n_2, \rng)$, we use $A \oplus B$ to denote the block-diagonal matrix
\[
A \oplus B := \begin{pmatrix} 
A & 0 \\
0 & B 
\end{pmatrix},
\]
of size $(m_1 + m_2) \times (n_1 + n_2)$.

Let $A=\begin{bmatrix}
	A_{11} & A_{12}\\
	0 & A_{22}
\end{bmatrix}\in\GL(n, \rng)$, where $A_{11}\in\M(n_1, \rng)$ and $A_{22}\in\M(n_2, \rng)$. Then $A_{11}$ and $A_{22}$ are also invertible, that is, $\det(A_{11})$ and $\det(A_{22})$ are units in $\rng$. To see this, let $\alpha=\det(A)$, $\beta=\det(A_{11})$, and $\gamma=\det(A_{22})$. So $\alpha=\beta\gamma$. As $\alpha$ is a unit, there exists $\alpha^{-1}$, such that $\alpha\cdot\alpha^{-1}=1$. So $(\beta\gamma)\alpha^{-1}=\beta(\gamma\alpha^{-1})=1$, so $\beta$ is also a unit.


\subsection{Rank notions of matrices over rings}\label{subsec:rank-notion}

Over fields, there are several equivalent definitions of the rank of a matrix, but some of these become inequivalent over more general rings. In this paper we will make use of two different notions of rank for rings over $R$. %
%
%We now recall several different equivalent notions of rank of matrices over fields, as they will not all be equivalent over more general rings. For a field $\F$ and $A\in\M(m\times n, \F)$, the following values of its rank all coincide:
%
%\begin{itemize}
%	\item The dimension of the linear span of rows (resp. columns) of $A$ as a subspace of $\F^m$ (resp. $\F^n$). 
%	\item The codimension ($n$, resp. $m$, minus dimension) of the left (resp. right) kernel of $A$ as a subspace of $\F^n$ (resp. $\F^m$). 
%	\item The smallest $k\in\N$ that allows for a decomposition of $A$ as $\sum_{i\in[k]}u_iv_i^t$, where $u_i\in \F^m$ and $v_i\in\F^n$. 
%	\item The largest $k\in\N$ such that $A$ has a $k\times k$ non-zero minor.
%	%\knote{do you mean largest?} 
%\end{itemize}
%
%Each of the notions above can be generalised to the ring setting, but they do not always give rise to rank notions, at least not easily. Let $A\in \M(m\times n, \rng)$. 
%\begin{itemize}
%	\item The $\rng$-span of rows (resp. columns) of $A$ as a submodule of $\rng^m$ (resp. $\rng^n$). 
%	\item The left (resp. right) kernel of $A$ as a submodule of $\rng^n$ (resp. $\rng^m$). 
%	\item The smallest $k\in\N$ that allows for a decomposition of $A$ as $\sum_{i\in[k]}u_iv_i^t$, where $u_i\in \rng^m$ and $v_i\in\rng^n$. 
%	\item The largest $k\in\N$ such that $A$ has a $k\times k$ non-zero minor.
%	%\knote{do you mean largest?}
%\end{itemize}
%
%To define ranks of modules over $\rng$ requires some caution. 
%
We refer the reader to \cite{Coh85} for free ideal rings over which ranks of modules could be defined. Here we only make use of the following rank notions. Let $A\in\M(m\times n, \rng)$.
\begin{enumerate}
	%	\item The linear independence row-rank of $A$, $\lirrk(A)$, is the maximum number of rows that are $\rng$-independent. Similarly the linear independence column-rank of $A$, $\licrk(A)$, can be defined. 
	
	%	Note that if $\lirrk(A)=m$, then the left kernel of $A$ is zero.
	
	\item The \emph{inner rank} of $A$, $\irk(A)$, is the smallest $r\in\N$ such that $A$ can be written as $\sum_{i\in[r]}u_iv_i^t$, where $u_i\in \rng^m$ and $v_i\in\rng^n$. 
	
	Or equivalently, the smallest $r\in \N$ such that $A=BC$ where $B\in\M(m\times r, \rng)$ and $C\in\M(r\times n, \rng)$. 
	
\item For $t\in \min\{m, n\}$, let $I_t(A)$ be the ideal of $\rng$ 
%\knote{'of $\rng$'} 
generated by $t\times t$ minors of $A$. For $t=0$, let $I_0(A)=\rng$. The \emph{McCoy rank} of $A$, $\mcrk(A)$, is the largest $r\in \{0, 1, \dots,  \min\{m, n\}\}$ such that the annihilator of $I_r(A)$ is the zero ideal \cite{McCoy-Rank}. 
\end{enumerate}

%For our purpose, we need a rank notion that is invariant under left and right multiplications of invertible matrices. By \cite{inner-rank-comp}, linear independence ranks do not satisfy this. 

To decide between inner and McCoy rank, the following example \cite[Section 4.2, Example 3]{inner-rank-comp} is useful. Consider $A=\diag(p, \dots, p)\in \M(n, \Z/p^{n+1}\Z)$. Then $\mcrk(A)=0$, as for $1\leq t\leq n$, $I_t(A)$ is generated by $p^t$, so its annihilator is generated by $p^{n+1-t}$. As seen below, $\irk(A)=n$, and this turns out to be useful for our purpose. 

However, McCoy rank is also useful in our context as it is closely related to the solvability of equations. Let $A\in \M(m\times n, \rng)$, and let $\vec{x}=(x_1, \dots, x_n)^t$. Then $Ax=\vec{0}$ has a non-zero solution over $\rng$ if and only if $\mcrk(A)<n$. In particular, when $A\in\M(n, \rng)$, $Ax=\vec{0}$ has %\knote{`non-trivially'?}  
non-zero solutions if and only if $\det(A)$ is a zero-divisor.

\subsection{Some basic facts of inner ranks}

Let $\rng$ be a commutative ring. Some basic properties of inner ranks are as follows \cite[pp. 248]{Coh85}. 
\begin{fact}\label{fact:inner-rank}
	Let $A \in \M(m \times n, \rng)$. Let $L\in\GL(m, \rng)$ and $R\in\GL(n, \rng)$. 
	\begin{enumerate}
		\item $0\leq \irk(A) \leq \min\{m,n\}$.
		\item $\irk(A\oplus B)\leq \irk(A)+\irk(B)$.
		\item $\irk(AB)\leq \min\{\irk(A), \irk(B)\}$. 
		\item $\irk(\begin{bmatrix}
			A & B
		\end{bmatrix})\geq \max\{\irk(A), \irk(B)\}$ and $\irk(\begin{bmatrix}
			A \\ B
		\end{bmatrix})\geq \max\{\irk(A), \irk(B)\}$.
		\item $\irk(A)=\irk(LAR)$.
	\end{enumerate}
\end{fact}

The inner rank is also equal %to the least $r$ such that $A$ can be written as a sum of $r$ matrices of the form $b \otimes c$ where $b$ is a column vector and $c$ is a row vector,  % \knote{there's no use of $r$ in this condition}, 
and 
to the least $r$ such that row-span of $A$ (that is, the $\rng$-submodule of $\rng^n$ generated by the rows of $A$) is contained in a submodule of $\rng^n$ that is generated by $r$ elements, and similarly for the column-span of $A$ \cite[Proposition~5.4.3]{Coh85}. The latter characterization cannot, in general, be strengthened to $\rho(A)=$ the minimum number of generators of the row-span of $A$ itself; see Example~\ref{ex:nonreduced} for a counterexample.

%\begin{observation} \label{obs:row-span}
%Let $A$ be an $m \times n$ matrix, and let $row(A)$ denote the $R$-submodule of $R^n$ generated by the rows of $A$, and similarly for $col(A) \leq R^m$. Then the inner rank of $A$ is equal to the smallest number $r$ such that there is an submodule of $R^n$ that contains $row(A)$ (resp., submodule of $R^m$ that contains $col(A)$) and is generated by $r$ elements.
%\end{observation}


%\begin{proof}
%Suppose $row(A)$ is generated as an $R$-submodule of $R^n$ by the row vectors $v_1, \dotsc, v_r \in R^n$. Then every row of $A$ can be written as $\sum_{j=1}^r x_j v_j$ for some $x_i \in R$. Let $B_{ij}$ be such that the $i$-th row of $A$ is the $R$-linear combination $\sum_{j=1}^r B_{ij} v_j$. If we let $C$ be the $r \times n$ matrix whose $i$-th row is $v_i$, then we have that $BC=A$, and those $\rho(A) \leq r$.
%
%Conversely, suppose $\rho(A) \leq r$. Then there are matrices $B\in\M(m\times r, \rng)$ and $C\in\M(r\times n, \rng)$ such that $A = BC$. Thus $row(A) \leq row(C)$, so $row(A)$ is contained in an $R$-submodule of $R^n$ that is generated by at most $r$ elements, namely, the rows of $C$.
%\end{proof}

\begin{example}[Inner rank can be strictly subadditive]\label{ex:Z6}
	For $\rng = \Z/6\Z$, consider the matrix $A = \begin{pmatrix} 2 & 0 \\ 0 & 3 \end{pmatrix}$. We have $A = \begin{pmatrix} 4 \\ 3 \end{pmatrix} \begin{pmatrix} 2 & 3 \end{pmatrix}$, so $\rho(A) \leq 1$ (in fact, it is equal to $1$, since only the zero matrix has inner rank zero). Thus, we do not in general have $\rho(B \oplus C) = \rho(B) + \rho(C)$. %\knote{This examples uses $A$ for two different things.}
\end{example}


We recall the following classical definitions, see, e.\,g., \cite[pp.~6--7]{Coh85}. 
\begin{enumerate}
	%	\item We say that $\rng$ has the \emph{invariant basis number} property (IBN), if the following holds: for any $A\in\M(m\times n, \rng)$ and $B\in\M(n\times m, \rng)$, if $AB=\idmat_m$ and $BA=\idmat_n$, then $m=n$.\footnote{The name comes from the fact that an equivalent condition is that $R^n \not\cong R^m$ for any $n \neq m$, so free modules have a well-defined size of their bases.}
	\item Let $\rng$ be a ring that is not necessarily commutative. We say that $\rng$ has the \emph{unbounded generating number} property (UGN), if the following holds: for any $A\in\M(m\times n, \rng)$ and $B\in\M(n\times m, \rng)$, if $AB=\idmat_m$, then $n\geq m$. Or equivalently, %\knote{`$\rng$ has UGN if'} 
	$\rng$ has UGN if
	$\rho(\idmat_n)=n$.\footnote{The name comes from the fact that an equivalent condition is that there is no constant upper bound on the number of generators of all $R$-modules.}
	%\knote{we show the equivalence mentioned in this footnote in Remark 3.2 so maybe we should add that to the footnote:  `; see Remark ...' or something}
	\item Let $\rng$ be a ring that is not necessarily commutative. We say that $\rng$ is \emph{stably finite} (sometimes called ``weakly finite''), if the following holds: for any $A, B\in \M(n, \rng)$, if $AB=\idmat_n$, then $BA=\idmat_n$.
\end{enumerate}
An observation due to P. M. Cohn \cite{Coh85} is that stable finiteness implies UGN. %, which in turn implies IBN. 
Stably finite rings include all %\knote{all?} 
commutative rings, Artinian rings, and Noetherian rings.

%Note that if a ring has UGN, then $\irk(\idmat_m)=m$. 
Following \cite[pp. 248]{Coh85}, we say that a matrix $A\in\M(n, \rng)$ is \emph{full} if $\irk(A)=n$, and it was also noted there that for a ring with UGN, every $A\in \GL(n, \rng)$ is full. 

\begin{lemma}\label{lem:upper-tri}
	Let $\rng$ be a commutative ring. Let $a, b, c\in \N$. Let $A\in \M(a, \rng)$ satisfy that $\mcrk(A)=a$. Suppose $(P, Q)\in\GL(a+b, \rng)\times\GL(a+c, \rng)$ satisfies that 
	\begin{equation}\label{eq:upper-tri}
		P\begin{bmatrix}
			A & 0 \\
			0 & 0_{b\times c}
		\end{bmatrix}Q^t=\begin{bmatrix}
			A & 0 \\
			0 & 0_{b\times c}
		\end{bmatrix}.
	\end{equation} Then $P=\begin{bmatrix}
		P_{11} & P_{12}\\
		0_{a\times b} & P_{22}
	\end{bmatrix}$, $Q=\begin{bmatrix}
		Q_{11} & Q_{12}\\
		0_{a\times c} & Q_{22}
	\end{bmatrix}$.
	%, where $P_{11}Q_{11}^t=\idmat_a$, $P_{22}\in\GL(b, \F)$, and $Q_{22}\in\GL(c, \F)$. 
\end{lemma}
\begin{proof}
	Let $P=\begin{bmatrix}
		P_{11} & P_{12}\\
		P_{21} & P_{22}
	\end{bmatrix}$, $Q=\begin{bmatrix}
		Q_{11} & Q_{12}\\
		Q_{21} & Q_{22}
	\end{bmatrix}$. Then Equation~\ref{eq:upper-tri} implies that 
	$$
	\begin{bmatrix}
		P_{11}AQ_{11}^t & P_{11}AQ_{21}^t\\
		P_{21}AQ_{11}^t & P_{21}AQ_{21}^t
	\end{bmatrix}
	=
	\begin{bmatrix}
		A & 0 \\
		0 & 0_{b\times c}.
	\end{bmatrix}
	$$
	%Note that\knote{`Note that' -> `This implies that'} 
	This implies that
	$P_{11}AQ_{11}^t=A$. As $\mcrk(A)=a$, $\det(A)=\det(P_{11}AQ_{11}^t)$ is a non-zero-divisor. It follows that each of $\det(P_{11})$, $\det(A)$, $\det(Q_{11}^t)$ are all non-zero-divisors. Therefore, $\mcrk(AQ_{11}^t)=a$, so by $P_{21}AQ_{11}^t=0$, we have $P_{21}=0$. To show that $Q_{21}=0$ follows the same line of argument.
	%	By $P_{11}Q_{11}^t=A$ and the stable finiteness of $\rng$ (since it is commutative), we have $Q_{11}^tP_{11}=A$. It follows that for any non-zero $v\in \rng^a$, $P_{11}v\neq 0$, as otherwise $Av=0$, contradicting that $\mcrk(A)=a$. So $P_{11}Q_{21}^t=0$ implies that $Q_{21}^t=0$. Similarly, for any non-zero $v\in\rng^a$, $v^tQ_{11}^t\neq 0$. So $P_{21}Q_{11}^t=0$ implies that $P_{21}=0$. 
\end{proof}

Inner ranks may have strange behaviours. For example, non-zero scalar matrices may not have full inner ranks.

\begin{example}[Inner ranks of (non-zero) scalar matrices over commutative Noetherian rings] \label{ex:nonreduced}
	Let $R = \F[x,y,z] / \langle x^2 - z, y^2 - z, xy \rangle$, where $\F$ is a field (though $\F=\Z$ would work as well). Since $x^2 = y^2 = z$ in $R$, but $xy=0$, we have $xz = xy^2 = 0$ and similarly $yz=0$, and $z^2 = x^2 y^2 = 0$. A little Gröbner basis calculation gives us that $\{x^2-z, y^2-z, xy, yz, xz, z^2\}$ is a Gröbner basis for the defining ideal here, and thus $\{1,x,y,z\}$ is an $\F$-basis of $R$ as a vector space. But we have $\begin{pmatrix} x \\ y \end{pmatrix} \begin{pmatrix} x & y \end{pmatrix} = z I_2$, hence $\rho(z I_2) = 1$, so a non-zero scalar matrix may not have full inner rank. This is a counterexample to scalar matrices of full inner rank for many classes of rings as, in addition to being commutative Noetherian, $R$ is local Artinian, every non-unit is a zero divisor, and every zero divisor is nilpotent. (To see that it is local, note that $x^4 = y^4 = z^2 = 0$, so geometrically this is a ``fat point at the origin;'' the unique maximal ideal is $\langle x,y \rangle = Span_\F\{x,y,z\}$.) We note that in this example, it is also the case that $row(z I_2)$ is not generated by 1 element despite $\rho(z I_2)=1$. (This example can be generalized to give a ring in which $\rho(z I_n) = n-1$ but that $\rho(r I_m) = m$ for all nonzero $r$ and all $m < n$.)
\end{example}

%\ynote{I moved the inner rank computation here. I could not find the LAA paper about SNF computation you mentioned...}
At some places, we may require that the inner rank of $A\in \M(\ell\times m, \rng)$ is efficiently computable. The inner rank of matrices over a principal ideal domain can be computed efficiently via the Smith Normal Form \cite{SmithNF,SmithNFJournal}. In \cite{inner-rank-comp}, it is shown that the inner rank of matrices over a finite commutative principal ideal ring can be computed in polynomial time.

%We will be interested in rings over which non-zero diagonal matrices are of full inner rank. This is true for $\Z$ (as it is a principal ideal domain) and $\Z/p^e\Z$ (as it is local).\ynote{This needs to be expanded.}

%
%\subsection{Efficiently computable commutative rings} 
%\jnote{Do we need this section at all? I think our results work with $R$-algorithms over any ring $R$... If that is correct, maybe we should just move the last paragraph about SNF and inner rank being computable in finite PIRs to the next section and get rid of this section}
%\ynote{Please check.}\knote{I suggest we move this section to the end of the Preliminaries, after the notation for matrices and and inner rank have been defined.} Let $\rng$ be a commutative ring. We impose some conditions to compute with $\rng$ efficiently. \ynote{Moved to here.}
%
%First, assume that the additive group of $\rng$, $\rng^+$, is finitely generated, so it is isomorphic to $A\oplus \Z^k$ where $A$ is a finite abelian group. We further assume that a set of generators for $A$ is given. 
%
%Second, the multiplication of $\rng$, that is a $\Z$-bilinear map $\phi: \rng^+\times\rng^+\to\rng^+$, is given by in the form of
%%\knote{suggest rephrase: `by' -> `in the form of'} 
%structure constants. 
%
%Third, we require that the inner rank of $A\in \M(\ell\times m, \rng)$ is efficiently computable. The inner rank of matrices over a principal ideal domain can be computed efficiently via the Smith Normal Form. In \cite{inner-rank-comp}, it is shown that the inner rank of matrices over a finite commutative principal ideal ring can be computed in polynomial time.
%\ynote{We should include an algorithm for computing inner ranks for matrices over finite rings and over $\Z$. Also how about the inner rank of matrices of linear forms?}

\subsection{Finite local principal ideal rings} \label{sec:PIR}

%\ynote{It should be good to include a general discussion on Smith Normal Form and the complexity of computing it. One reference I found is \cite{SmithNF}.}\jnote{I believe that got published in LAA in 1998 so we should cite the journal version. One potential drawback is I have not checked whether the algorithms there actually work over arbitrary PIRs, or whether some things are specific to PIDs or even to $\Z$ itself. Either way we should still cite it. If we find that the algorithms there do not work for PIRs, we should cite some standard textbook that reviews the naive Gaussian-elimination-style algorithm. I believe this is covered in Artin's Algebra textbook, when he does the Fundamental Thm of Abelian Groups. I will try to dig through all of this...}

Let $\rng$ be a finite local PIR. We will compute with modules over finite local PIRs and their tensor products in Section~\ref{subsec:five-action-PIR} and Section~\ref{subsec:tuple-isometry}. Let us first record some examples and notations.
\begin{itemize}
	\item 
	Examples of such rings include $\Z/p^e \Z$, $\F[x] / x^n$ where $\F$ is a finite field, and Galois rings.%\ynote{I removed the odd characteristic conditions because here we do not assume odd characteristic.}
	
%	with $p$ odd (the crucial case for the application to \algprobm{Group Isomorphism} we need), $\F[x] / x^n$ where $\F$ is a field of odd characteristic, %the $p$-adic integers $\Z_p$ for $p$ odd, 	and Galois rings of odd characteristic.
	
	\item The unique maximal ideal of $\rng$ is denoted $\m$.
	
	\item We use $p$ to denote the characteristic of $\rng/\m$.
	%; by assumption $p$ is an odd prime.
	
	\item We use $e$ to denote the index of nilpotency of $\m$: $\m^e = 0$ but $\m^{e-1} \neq 0$. (Since $R$ is finite and local, $\m$ is necessarily nilpotent.)
	
	\item When working over a ``general'' finite local PIR, we treat basic arithmetic operations in $R$ as unit cost (or, equivalently, we work in the BSS model over $R$ \cite{BSS}). We also assume that division in $R$ is unit cost: given $a,b \in R$, decide if $a$ divides $b$, and if so, output some $c$ such that $ac = b$. We also assume that we have access to $t$, a generator of the maximal ideal. We refer to algorithms in this model as $R$-algorithms. 
	
	\item We also assume that elements $x \in R/\m^i$ are represented computationally by elements $\overline{x} \in R$ such that $\overline{x} \equiv x \pmod{\m^i}$, and the algorithm simply keeps track of when it is working modulo $\m^i$. For example, to test whether $\overline{x}, \overline{y}$ are equivalent modulo $\m^i$, an $R$-algorithm can simply check whether $\overline{x} - \overline{y}$ is divisible by $t^i$. 
		
\end{itemize}

Smith Normal Form is efficiently computable by an $\rng$-algorithm, by the standard ``Gaussian-elimination-style'' algorithm, for matrices over a finite local PIR $\rng$.\footnote{Lacking a convenient reference that does not assume either the integers or at least a domain, we outline it here. Since a local commutative PIR is a quotient of a PID \cite{hungerford}, the algorithm does the same operations as though one is over a PID, keeping track of which elements are equal or zero in the quotient PIR. In outline: standard row operations can make the (1,1) entry into the GCD of the entries in the first column, and then multiples of the first row can be added to the remaining rows to zero out the rest of the first column. This procedure can be repeated on the second column, to make the (2,2) entry the GCD of the new entries of the second column at or below the second row. And so on. The result is an upper triangular matrix. Now we repeat the analogous procedure but using column operations, resulting in a diagonal matrix.} In many cases of interest, including $\rng = \Z/p^e \Z$, %$R$ is finite and 
ring operations in $\rng$ are computable in $\poly(\log|\rng|)$ time on an ordinary Turing machine. With a little more work this is probably also true over arbitrary finite local PIRs, but we will not need this in our paper.

An $\rng$-algebra $A$ is a unital ring $A$ together with an injective ring homomorphism $\iota \colon \rng \to A$, such that the image of $\iota$ (which includes the multiplicative identity $1 \in A$) is contained in the center of $A$. Because of this, the underlying additive group of $A$ is an $\rng$-module. A \emph{$*$-algebra over $\rng$} is an $\rng$-algebra $A$ together with an anti-automorphism $* \colon A \to A$ of order at most 2, that is $(a+b)^* = a^* + b^*$, $(ab)^* = b^* a^*$, and $(a^*)^* = a$ for all $a \in A$. 

Given an $R$-module $M$, its endomorphism ring $\End(M)$ consists of all $\rng$-module homomorphisms $M \to M$. In this ring, addition is pointwise---$(f+g)(x) = f(x) + g(x)$---and multiplication is given by the composition of homomorphisms. A commutative ring has all its modules direct sums of cyclic modules (this condition is sometimes called being a ``Köthe ring'') if and only if it is an Artinian PIR \cite{CohenKaplansky, Kothe}. In particular, this holds for our ring $\rng$ by assumption (since finite implies Artinian), and thus every $\rng$-module has the form $M \cong \bigoplus_i R/\m^{e_i}$ for some integers $1 \leq e_i \leq e$. If $M$ is finitely generated, there are only finitely many such summands, say $n$ of them. 

The following theorem is the natural generalization of Hillar \& Rhea's description of the endomorphism ring and automorphism group of Abelian $p$-groups \cite{HillarRhea} (originally proved by Ranum 100 years earlier \cite{Ranum}) to the endomorphism ring and automorphism group of $\rng$-modules, when $\rng$ is a commutative local PIR. 

\begin{theorem} \label{thm:HillarRheaPIR}
	Let $\rng$ be a commutative local PIR with maximal ideal $\m = \langle t \rangle$. For $M = \bigoplus_{i=1}^n \rng/\m^{e_i}$, $\End(M)$ is isomorphic to the ring of $n \times n$ matrices whose $(i,j)$ entry is an element of $\Hom(\rng/\m^{e_j}, \rng/\m^{e_i})$ for all $i,j \in \{1,\dotsc,n\}$. 
	
	Furthermore, this is isomorphic to the ring of $n \times n$ matrices whose $(i,j)$ entry is an element of $\rng/\m^{e_i}$, that is divisible by $t^{e_i - e_j}$ whenever $e_i > e_j$. Such a matrix is invertible if and only if it is invertible when every entry is taken modulo $\m$.
\end{theorem}

As the proof is nearly identical to that of Hillar and Rhea in the case of $\rng = \Z/p^e \Z$ \cite{HillarRhea}, we place it in Appendix~\ref{app:similar}.

For an $\rng$-module $M$, a \emph{basis} is a generating set $m_1, \dotsc, m_n$ with the property that $\langle m_i \rangle \cap \langle m_j : j \neq i \rangle = 0$ for all $i=1,\dotsc,n$, where we use $\langle$ angle brackets $\rangle$ to denote the generated $\rng$-module. In particular, if $\{m_1,\dotsc,m_n\}$ is a basis of $M$, we have $M = \bigoplus_{i=1}^n \langle m_i \rangle$, and $M$ is a direct sum of cyclic $\rng$-modules. Over general commutative rings, not every finitely generated module has such a basis (e.g., the ideal $\langle x, y\rangle \subseteq \F[x,y]$, as an $\F[x,y]$-module), but because modules over commutative PIRs are direct sums of cyclic modules, every module over our ring $\rng$ has a basis. %(In fact, these are equivalent conditions: a commutative ring $\rng$ is a PIR if and only if all of its modules are direct sums of cyclic modules \cite{TODO}.)
